\chapter*{Annexes}
\label{annexes}
\phantomsection
\addcontentsline{toc}{chapter}{Annexes}

\section*{Annexe A : Tableau complet des références composants}
\begin{tableau}[H]
    \centering
    \begin{tabular}{|l|l|l|}
        \hline
        \textbf{Catégorie} & \textbf{Référence} & \textbf{Fonction} \\
        \hline
        Diodes & Vishay V35PW22HM3 & Redressement AC/DC (PD3) \\
        \hline
        Microcontrôleur & PIC18F25K80 & Supervision et commande CAN \\
        \hline
        Transceiver CAN & TJA1050 & Interface physique du bus CAN \\
        \hline
    \end{tabular}
    \caption{Extrait de la nomenclature (BOM) détaillée.}
\end{tableau}

\section*{Annexe B : Détails de calculs longs}
Les détails du dimensionnement des composants inductifs et capacitifs pour le convertisseur SEPIC et l'étude thermique complète des dissipateurs ($R_{thHA} = \SI{13.92}{\celsius\per\watt}$) sont regroupés ici. Ces étapes n'ont pas été détaillées dans le corps du texte pour des raisons de concision.

\section*{Annexe C : Trames CAN détaillées}
Le réseau CAN fonctionne à une vitesse de \SI{500}{\kilo\bit\per\second}. Le format des trames standard (CAN 2.0A) est exploité avec les identifiants (ID) suivants:
\begin{itemize}
    \item \texttt{ID 0x100} : Mesure de tension principale.
    \item \texttt{ID 0x101} : Mesure de courant.
    \item \texttt{ID 0x200} : Commande des relais (éclairage, USB).
\end{itemize}

\section*{Annexe D : Captures de simulation complémentaires}
Cette annexe recueille les relevés chronogrammes PSIM (tension, courant) pour la plage de charge complète (10\% à 100\%), qui n'ont pas été intégrés dans la section d'étude de la conversion afin de limiter la densité graphique du rapport.

\section*{Annexe E : Contraintes fabrication PCB}
Les circuits imprimés respectent des règles de routage strictes pour minimiser les perturbations EMI des composants de puissance envers la logique (largeur de piste, isolation, vias de dissipation thermique sous les composants CMS de puissance).

